\subsection{Finite-Dimensional Topological Vector Spaces}

% **Role:** topological addendum to finite-dimensional linear structure.
%
% Topological vector spaces, standard topology on R^n, natural topology on a finite-dimensional real vector space.

Finite-dimensional real vector spaces carry a canonical topology: choose coordinates, use the standard Euclidean topology on $\R^n$, and then forget that coordinates were chosen. All coordinate changes are homeomorphisms.

\begin{definition} \label{def:topological-vector-space} \hlt{Topological Vector Space}\\
Let $\mathcal{V}$ be a real vector space and let $\mathcal{T}$ be a topology on its underlying set.\\
The pair $(\mathcal{V},\mathcal{T})$ is a \hlt{topological vector space} if addition $A:\mathcal{V}\times\mathcal{V}\to\mathcal{V}$, $A(\mathbf{u},\mathbf{v})=\mathbf{u}+\mathbf{v}$, and scalar multiplication $M:\R\times\mathcal{V}\to\mathcal{V}$, $M(a,\mathbf{v})=a\mathbf{v}$ are continuous, where products carry the product topology.
\end{definition}

\begin{example} \label{ex:rn-standard-topological-vector-space} \hlt{The Standard Topology on $\R^n$}\\
With the Euclidean topology, $\R^n$ is a topological vector space: the maps $(\mathbf{x},\mathbf{y})\longmapsto \mathbf{x}+\mathbf{y}$, $(a,\mathbf{x})\longmapsto a\mathbf{x}$ have polynomial coordinate functions $x_i+y_i$ and $ax_i$, hence are continuous.
\end{example}

\begin{proposition} \label{prop:real-matrix-maps-continuous-homeomorphisms} \hlt{Real Matrix Maps Are Continuous}\\
Let $A=[a_{j,k}]\in\operatorname{Mat}(m,n,\R)$, and let $T_A:\R^n\to\R^m$ be the induced map $T_A(\mathbf{x})=A\mathbf{x}$. Then $T_A$ is continuous for the Euclidean topologies. If $m=n$ and $A$ is invertible, then $T_A$ is a homeomorphism. In particular, every linear functional on $\R^n$ is continuous.
\begin{proof}
For $\mathbf{x}=(x_1,\ldots,x_n)^t$, the $j$th coordinate of $T_A(\mathbf{x})$ is $\sum\limits_{k=1}^n a_{j,k}x_k$, hence is continuous. Thus $T_A$ is continuous. If $A$ is invertible, then $(T_A)^{-1}=T_{A^{-1}}$, which is continuous by the first part. Finally, every $f\in(\R^n)'$ equals $T_A$ for a unique $A\in\operatorname{Mat}(1,n,\R)$, so $f$ is continuous.
\hlt{Structural Commentary}: Matrix maps on Euclidean coordinate spaces are continuous coordinate by coordinate.\\
\hlt{Key Insight}: Invertible real matrices give homeomorphisms of Euclidean space.
\end{proof}
\end{proposition}

\begin{definition} \label{def:natural-topology-finite-dimensional-real-vector-space} \hlt{Natural Topology}\\
Let $\mathcal{V}$ be an $n$-dimensional real vector space, let $\beta$ be a basis, and let $[\cdot]_{\beta}:\mathcal{V}\to\R^n$ be the coordinate isomorphism.\\
The \hlt{natural topology} on $\mathcal{V}$ is $\mathcal{T}_{\beta} \defeq \{[\cdot]_{\beta}^{-1}(O): O\subseteq\R^n$ is open$\}$.\\
Equivalently, $U\subseteq\mathcal{V}$ is open if and only if $[\cdot]_{\beta}(U)$ is open in $\R^n$.
\end{definition}

\begin{proposition} \label{prop:natural-topology-independent-of-basis} \hlt{Natural Topology Is Basis-Independent}\\
If $\beta$ and $\gamma$ are bases of $\mathcal{V}$, then $\mathcal{T}_{\beta}=\mathcal{T}_{\gamma}$.
\begin{proof}
Set $H_{\gamma\leftarrow\beta}\defeq[\cdot]_{\gamma}\circ[\cdot]_{\beta}^{-1}:\R^n\to\R^n$.\\
Then $H_{\gamma\leftarrow\beta}\mathbf{x}=C_{\gamma\leftarrow\beta}\mathbf{x}$. Since $C_{\gamma\leftarrow\beta}$ is invertible with inverse $C_{\beta\leftarrow\gamma}$, Proposition \ref{prop:real-matrix-maps-continuous-homeomorphisms} shows that $H_{\gamma\leftarrow\beta}$ is a homeomorphism with inverse $H_{\beta\leftarrow\gamma}$.\\
Hence, for $U\subseteq\mathcal{V}$, $U\in\mathcal{T}_{\beta}$ $\iff$ $[\cdot]_{\beta}(U)$ is open $\iff$ $H_{\gamma\leftarrow\beta}([\cdot]_{\beta}(U))=[\cdot]_{\gamma}(U)$ is open $\iff$ $U\in\mathcal{T}_{\gamma}$.
\hlt{Structural Commentary}: Basis changes are homeomorphisms.\\
\hlt{Key Insight}: The natural topology is carried by the vector-space structure, not by a preferred basis.
\end{proof}
\end{proposition}

\begin{theorem} \label{thm:natural-topology-characterization} \hlt{Characterization of the Natural Topology}\\
Let $\dim_{\R}\mathcal{V}<\infty$. There is a unique topology $\mathcal{T}$ on $\mathcal{V}$ such that:
\begin{enumerate}[label=(\roman*)]
\setlength{\itemsep}{0pt}
\item $(\mathcal{V},\mathcal{T})$ is a topological vector space;
\item $\mathcal{V}'\subseteq C((\mathcal{V},\mathcal{T}),\R)$.
\end{enumerate}
Moreover, $\mathcal{T}=\mathcal{T}_{\beta}$ for every basis $\beta$.
\begin{proof}
Fix an arbitrary basis $\beta=(\mathbf{v}_1,\ldots,\mathbf{v}_n)$.
\begin{enumerate}[label=(\roman*)]
\setlength{\itemsep}{0pt}
\item \hlt{Existence}: put $\mathcal{T}=\mathcal{T}_{\beta}$. Since $[\cdot]_{\beta}:(\mathcal{V},\mathcal{T}_{\beta})\to\R^n$ is a homeomorphism, the vector operations on $\mathcal{V}$ are transported from the continuous vector operations on $\R^n$.\\
Also, for $\varphi\in\mathcal{V}'$, $\varphi\circ[\cdot]_{\beta}^{-1}\in(\R^n)'$, hence is continuous by Proposition \ref{prop:real-matrix-maps-continuous-homeomorphisms}. Thus (i)--(ii) hold.
\item \hlt{Uniqueness}: suppose $\mathcal{T}$ satisfies (i)--(ii).\\
For $1\leq i\leq n$, the coordinate functional $\pi_i\circ[\cdot]_{\beta}\in\mathcal{V}'$ is continuous, so $[\cdot]_{\beta}:(\mathcal{V},\mathcal{T})\to\R^n$ is continuous.\\
Its inverse $\Psi(a_1,\ldots,a_n)=a_1\mathbf{v}_1+\cdots+a_n\mathbf{v}_n$ is continuous by (i), since it is built from scalar multiplication and finite sums. Hence $[\cdot]_{\beta}$ is a homeomorphism, so $\mathcal{T}=\mathcal{T}_{\beta}$.
\end{enumerate}
\hlt{Structural Commentary}: Functionals give coordinates; operations give inverse coordinates.\\
\hlt{Key Insight}: Functionals and vector operations determine finite-dimensional coordinates topologically.
\end{proof}
\end{theorem}

\begin{corollary} \label{cor:finite-dimensional-linear-maps-continuous} \hlt{Finite-Dimensional Linear Maps Are Continuous}\\
Let $\dim_{\R}\mathcal{V},\dim_{\R}\mathcal{W}<\infty$, with natural topologies. Then $\LL(\mathcal{V},\mathcal{W})\subseteq C(\mathcal{V},\mathcal{W})$.\\
If $T\in\LL(\mathcal{V},\mathcal{W})$ is an isomorphism, then $T$ is a homeomorphism.
\begin{proof}
Choose bases $\beta$ of $\mathcal{V}$ and $\gamma$ of $\mathcal{W}$. For $T\in\LL(\mathcal{V},\mathcal{W})$, $[\cdot]_{\gamma}\circ T\circ[\cdot]_{\beta}^{-1}(\mathbf{x}) =[T]_{\gamma\leftarrow\beta}\mathbf{x}$, which is continuous by Proposition \ref{prop:real-matrix-maps-continuous-homeomorphisms}. Since $[\cdot]_{\beta}$ and $[\cdot]_{\gamma}$ are homeomorphisms, $T$ is continuous.\\
If $T$ is an isomorphism, apply the same argument to $T^{-1}$.\\
\hlt{Structural Commentary}: Finite-dimensional linear maps are just matrix maps in natural coordinates.\\
\hlt{Key Insight}: A basis turns topological questions about linear maps into continuity of matrix multiplication.
\end{proof}
\end{corollary}

\begin{perspective} \label{per:finite-dimensional-topology-is-rigid} \hlt{Finite-Dimensional Topology Is Rigid}\\
For finite-dimensional real vector spaces, topology is algebraically forced.\\
Note that $\mathcal{T}_{\beta}=\mathcal{T}_{\gamma}$; $\LL(\mathcal{V},\mathcal{W})\subseteq C(\mathcal{V},\mathcal{W})$; $T\in\LL(\mathcal{V},\mathcal{W})$ bijective $\Rightarrow$ $T$ is a homeomorphism.\\
Infinite-dimensional topological vector spaces need not be rigid in this sense.
\end{perspective}
